\documentclass[11pt,a4paper]{article}
\usepackage[utf8]{inputenc}
\usepackage{amsmath,amssymb}
\usepackage{geometry}
\geometry{margin=2cm}
\usepackage{array}
\usepackage{booktabs}
\usepackage{longtable}
\usepackage{enumitem}
\usepackage{xcolor}
\usepackage{natbib}
\usepackage[hidelinks]{hyperref}
\usepackage{parskip}

\sloppy
\definecolor{bettergrey}{gray}{0.95}
\definecolor{bettergreen}{RGB}{29, 191, 47}
\definecolor{betterblue}{RGB}{33, 65, 196}
\newcommand{\greenText}[1]{\textcolor{bettergreen}{#1}}
\newcommand{\blueText}[1]{\textcolor{betterblue}{#1}}
\newcommand{\field}[1]{\nolinkurl{#1}}
\newcommand{\importantfield}[1]{\textcolor{betterblue}{\nolinkurl{#1}}}
\newcommand{\icd}[1]{\texttt{#1}}
\renewcommand{\labelitemi}{$\vcenter{\hbox{\tiny$\bullet$}}$}
\renewcommand{\labelitemii}{$\scriptscriptstyle\square$}
\newcommand{\mr}{\specialrule{0.4pt}{4pt}{4pt}}
\newcommand{\tr}{\specialrule{0.8pt}{4pt}{6pt}}
\newcommand{\phase}[2]{%
  \subsection*{\textcolor{betterblue}{#1}: #2}%
  \addcontentsline{toc}{subsection}{#1: #2}%
}

\begin{document}

\begin{center}
	{\huge \textbf{Feasibility Analysis}}\\[1em]
	{\Large MIMIC-IV QA Benchmark for Coverage vs. Dispersion}\\[4em]
\end{center}

\tableofcontents
\newpage

%% ============================================================
\section{Introduction}
\label{sec:introduction}

The experiment's goal is that \emph{coverage} (facility-location) outperforms \emph{dispersion} (MMR, gMMR, FPS) when the candidate pool of retrieved chunks has a multi-cluster structure. Other QA datasets exist (HotpotQA, 2WikiMultiHopQA, MuSiQue), but these benchmarks have fixed candidate pools that we cannot control.

\paragraph{Ideal Cluster Structure for coverage}
\begin{enumerate}[leftmargin=*]
	\item \textbf{Redundant central cluster} - dominant topic, many similar chunks. Top-$k$ over-samples from here.  Coverage naturally avoids redundancy.
	\item \textbf{Complementary cluster(s)} - other relevant aspects of the query. Coverage ensures at least one representative from each.
	\item \textbf{Niche cluster} - small but informative, not an outlier. Coverage includes it because facility-location rewards covering every region of $D$.
	\item \textbf{Outliers} - informative or noisy.  Dispersion favors them (maximizes spread); coverage ignores them unless they represent a dense region.
\end{enumerate}

Proposal: build a \textbf{custom QA benchmark from MIMIC-IV} where we can (a)~verify that candidate pools have the desired cluster structure, (b)~control pool composition, and (c)~design multi-aspect queries that require many sources.

%% ============================================================
\part{MIMIC-IV Dataset}
\label{sec:data}

The database is split across two datasets (important CSVs highlighed in \textcolor{betterblue}{blue}):
\begin{itemize}[leftmargin=*]
	\item \textbf{MIMIC-IV} $\Rightarrow$ tabular clinical data recording every diagnosis, procedure, and medication for each hospitalization. This can give us \textbf{ground truth} for evaluation. There are two submodules:

	      \begin{itemize}[itemsep=1pt, parsep=0pt]
		      \item \textbf{hosp}: patient tracking (\importantfield{patients}, \importantfield{admissions}, \field{transfers}), diagnoses (\importantfield{diagnoses\_icd}, \importantfield{d\_icd\_diagnoses}, \field{procedures\_icd}, \field{drgcodes}), labs (\field{labevents}), medications (\field{prescriptions}, \field{emar}).
		      \item \textbf{icu}: bedside data (\field{chartevents}, \field{inputevents}, \field{outputevents}, \field{procedureevents}).
	      \end{itemize}

	\item \textbf{MIMIC-IV-Note} $\Rightarrow$ free-text clinical notes
	      \begin{itemize}[itemsep=1pt, parsep=0pt]
		      \item \importantfield{discharge}: discharge summaries. Long, narrative clinical documents written by doctors when a patient leaves the hospital.
		      \item \field{radiology}: radiology reports.
	      \end{itemize}
\end{itemize}

\bigskip\hrule\bigskip
For our purposes, the critical tables are:
\begin{itemize}[itemsep=2pt, parsep=1pt]
	\item \importantfield{discharge} (in MIMIC-IV-Note): $\sim$331K discharge summaries.
	\item \importantfield{diagnoses\_icd}: join table linking ICD-9/ICD-10 (International Classification of Diseases) codes to subject hospitalizations, with approximate importance ordering (\field{seq\_num}). A typical admission has 5-20 ICD codes, often spanning multiple body systems. Definitions for the codes themselves are in the \importantfield{d\_icd\_diagnoses} table
	\item \importantfield{patients}: gender, \field{anchor\_age}, date of death.
	\item \importantfield{admissions}: admission type, insurance, language, marital status, race, discharge location.
\end{itemize}

\vspace{3mm}

All tables link via \field{subject\_id} (patient) and \field{hadm\_id} (hospitalization). The dataset has 331,793 discharge summaries from 145,914 patients. All protected health information (names, dates, locations) has been removed or replaced with placeholders (\_\_\_).

Some analytics about demographics:

\begin{center}
	\begin{tabular}{lrr}
		                                     & \textbf{Hospital admissions}   & \textbf{ICU admissions} \\
		\mr
		Number of stays                      & 546,028                        & 94,458                  \\
		Unique patients                      & 223,452                        & 65,366                  \\
		Age, mean (SD)                       & 59.2 (19.1)                    & 64.7 (16.9)             \\
		Female Administrative Gender, n (\%) & 284,097 (52.0)                 & 37,714 (44.2)           \\
		Hospital length of stay, mean (SD)   & 4.8 (7.2)                      & 3.6 (5.4)               \\
		In-hospital mortality, n (\%)        & 11,801 (2.2)                   & 9,475 (11.1)            \\
		One year mortality, n (\%)           & 64,934 (11.9)                  & 14,951 (17.5)           \\
		\mr
		Data period                          & \multicolumn{2}{c}{2008--2022}                           \\
	\end{tabular}
\end{center}

\section{MIMIC-IV-Note}

The notes are of two kinds: \textbf{Discharge Summaries} and \textbf{Radiology Reports}.

\subsection{Discharge Summaries}
\label{sec:note-structure}

The structure of the notes seems consistent, with standardized sections:

\bigskip
\begin{center}
	\small
	\begin{tabular}{p{4cm} p{6.5cm}p{3.8cm}}
		\textbf{Section}                          & \textbf{Content}                                   & \textbf{QA Value}                                   \\
		\tr
		\blueText{\textbf{Brief Hospital Course}} & \blueText{\textbf{Clinical reasoning per problem}} & \blueText{\textbf{Highest}}                         \\
		History of Present Illness                & Symptom narrative, temporal progression            & High                                                \\
		Pertinent Results                         & Labs, imaging, micro interpretations               & High                                                \\
		Discharge Diagnosis                       & Diagnostic labels                                  & High (linkable to ICD)                              \\
		\mr
		Chief Complaint                           & 1--2 lines, reason for admission                   & \textbf{High} only as contextual embedding metadata \\
		Past Medical History                      & Comorbidities list                                 & Medium                                              \\
		Physical Exam                             & Admission + discharge findings                     & Medium                                              \\
		Discharge Medications                     & Treatment protocols with dosing                    & Medium                                              \\
		\mr
		Major Surgical or Invasive Procedure      & Procedures performed during stay                   & Low                                                 \\
		Allergies                                 & -                                                  & Low                                                 \\
		Family History                            & -                                                  & Low                                                 \\
		Medications on Admission                  & -                                                  & Low                                                 \\
		Discharge Disposition                     & Destination (home, SNF, rehab, \ldots)             & Low                                                 \\
		Discharge Condition                       & Mental status, consciousness, activity             & Low                                                 \\
		Discharge Instructions                    & Patient-facing plain-language summary              & Low                                                 \\
		Service                                   & Medicine, Surgery, \ldots                          & None                                                \\
		Attending                                 & Attending physician                                & None (deidentified)                                 \\
		Social History                            & Lifestyle, substance use                           & None (deidentified)                                 \\
		Followup Instructions                     & Outpatient follow-up appointments                  & None (deidentified)                                 \\
	\end{tabular}
\end{center}
\bigskip

\paragraph{Brief Hospital Course.}
This section uses \texttt{\#} markers to delimit individual clinical problems (e.g., \texttt{\# Stroke}, \texttt{\# Hypertension}, \texttt{\# DVT}).  Each sub-section contains \emph{clinical reasoning}: differential diagnoses, treatment rationale, management decisions with justifications.  This is where generalizable medical knowledge is embedded within patient-specific narratives.

However, the \texttt{\#}-header format is not universal. There are two patterns:

\begin{enumerate}[leftmargin=*]
	\item \textbf{Complex medicine admissions} (the interesting cases):
	      the Hospital Course is organized as a problem list with
	      \texttt{\#} headers, one block per active clinical issue.
	      Example: a cirrhosis patient's Hospital Course has
	      \texttt{\# ASCITES}, \texttt{\# HEPATIC ENCEPHALOPATHY},
	      \texttt{\# HYPONATREMIA}, \texttt{\# HIV}, \texttt{\# COPD}
	      (5--8 blocks, each 50--200 tokens).  Sometimes there are
	      additional sub-headers like \texttt{CHRONIC ISSUES} or
	      \texttt{ACUTE/ACTIVE ISSUES} grouping the problems.
	      The formatting is inconsistent: \texttt{\# Ascites},
	      \texttt{\# ASCITES}, \texttt{\#NSTEMI} (with or without space
	      after \texttt{\#}, mixed casing).

	\item \textbf{Simple or surgical admissions} (single-problem cases):
	      the Hospital Course is a single narrative paragraph with \emph{no} \texttt{\#} headers.  This includes orthopedic (hip fracture), urological (nephrectomy), and straightforward medicine cases (dysphagia workup).  These are analogous to single-hop queries. There is only one information peak, so coverage should add nothing.
\end{enumerate}

\textbf{The \texttt{TRANSITIONAL ISSUES} block.}  This (sometimes misspelled as TRANISTIONAL ISSUES) appears at the
end of many Hospital Courses and contains follow-up action items
(``recheck electrolytes'', ``follow up with hepatology''), not clinical
narratives.  These are task-oriented and do not describe clinical problems.

\paragraph{Deidentification artifacts.}
The \texttt{\_\_\_} placeholder is everywhere in dates, names, ages, and locations.  Brief Hospital Course and Physical Exam can contain few of them but largely are usable. Ages are recoverable from \field{anchor\_age} in the \field{patients} table.

%% ============================================================
\subsection{Radiology Reports}
\label{sec:radiology}

The \field{radiology} table in MIMIC-IV-Note contains $\sim$2.3M radiologist reports spanning x-ray, CT, MRI, and ultrasound.  Reports follow a structured format with sections: \textbf{Examination}, \textbf{Indication}, \textbf{Technique}, \textbf{Comparison}, \textbf{Findings}, and \textbf{Impression}.  The associated \field{radiology\_detail} table provides exam codes and exam names (e.g., \texttt{CT HEAD W/O CONTRAST}, \texttt{LIVER OR GALLBLADDER US}) that can be used for filtering.

\paragraph{Content quality is dual:}
\begin{itemize}[itemsep=2pt, parsep=1pt]
	\item \textbf{Formulaic reports} (routine chest x-rays, follow-up films): short, low information content.  Typical impression: ``\emph{No acute cardiopulmonary abnormality.}''
	\item \textbf{Detailed diagnostic reports} (CT head, CT abdomen with contrast, MRI brain, CT angiography): substantial findings sections
	      with organ-by-organ descriptions, measurements, and clinical correlation.  Example from a CT head: ``\emph{Subcortical and periventricular white matter hypodensities consistent with chronic small vessel ischemic disease.  No evidence of acute intracranial hemorrhage, edema, mass effect, or large territorial infarction.}''
\end{itemize}

\paragraph{Potential value for coverage vs.\ dispersion.}
For imaging-dependent conditions (stroke, PE, pneumonia), radiology reports provide a \textbf{different semantic perspective} from discharge notes: the discharge note describes clinical management (``tPA administered within 3h window''), while the radiology report describes anatomical findings (``acute hypodensity in L MCA territory, no midline shift''). This seems good for coverage.

\paragraph{Redundancy with discharge notes.}  The ``Pertinent Results'' section of discharge summaries often summarizes key imaging findings.
Radiology reports may not add unique facets beyond what is already captured there.

I think Discharge Notes are better as a starting point and this can be used later if needed.
\bigskip

%% ============================================================
\subsection{Considerations}

\subsubsection{Patient-specific vs.\ Generalizable data}
\label{sec:tension}

Discharge notes are patient-specific narratives, not medical textbooks.  A query like ``\emph{what are the general symptoms of a stroke?}'' has no ground truth in this corpus: no individual chunk states textbook symptomatology.  Embedding similarity will match chunks that \emph{mention} stroke, but the question is unanswerable from patient notes alone.

The queries should ask about \textbf{clinical management patterns}: things that are described in discharge notes (the Brief Hospital Course says what was done and why).  A more suitable query:
\begin{quote}
	\small
	``\emph{How is pneumonia managed in an elderly patient with chronic kidney failure?}''
\end{quote}

This works because: (a)~management decisions (antibiotic choice, fluid balance, renal dosing) are explicitly documented in discharge notes; (b)~answering requires synthesizing across multiple patients: some with pneumonia alone, some with pneumonia + CKD, providing both the baseline and the modifier; (c)~ground truth is derivable: chunks from admissions carrying both pneumonia and CKD ICD codes, plus chunks from pneumonia-only admissions for contrast.  Each chunk contributes one patient's management; the gold set is the collection of chunks whose union covers all required facets.

This query style aligns with the coverage hypothesis.  Top-$k$ will over-retrieve the most common pneumonia presentation (central cluster). \textbf{Coverage} ensures representation from CKD-modified management, elderly-specific complications (delirium, frailty), and atypical presentations.

%% ============================================================
\section{MIMIC-IV Structured Data}
\label{sec:structured-data-detail}

\subsection{Hosp Module}
\label{sec:hosp}

The \textbf{hosp} module captures hospital-wide data from admission to discharge.
Key table groups (tables \textcolor{betterblue}{highlighted in blue} are useful for us):

\paragraph{Patient tracking.}
Three tables describe demographics and movement:
\begin{itemize}[itemsep=2pt, parsep=1pt]
	\item \importantfield{patients}: gender, \field{anchor\_age} (deidentified age), date of death.
	      Ages are shifted but intervals are preserved, so age \emph{at admission} is recoverable.
	\item \importantfield{admissions}: admission/discharge times, admission type (emergency, elective, urgent),
	      insurance, language, marital status, race, discharge location, in-hospital mortality flag.
	\item \field{transfers}: ward-level ADT (admission--discharge--transfer) records tracking
	      physical location throughout the stay.
\end{itemize}

\paragraph{Billing and diagnoses.}
\begin{itemize}[itemsep=2pt, parsep=1pt]
	\item \importantfield{diagnoses\_icd}: up to 39 ICD-9-CM / ICD-10-CM codes per admission with
	      approximate importance ordering (\field{seq\_num}).  Definitions in \importantfield{d\_icd\_diagnoses}.
	      This is our primary source for condition-level labels and query scaffolding.
	\item \field{procedures\_icd}: billed procedures (ICD-9-PCS / ICD-10-PCS).
	      Definitions in \field{d\_icd\_procedures}.
	\item \field{drgcodes}: Diagnosis Related Group codes (AP-DRG and MS-DRG) assigning an
	      overall cost category to each hospitalization.
	\item \field{hcpcsevents}: HCPCS billing codes for provided services (e.g.\ mechanical
	      ventilation, ICU care).  Definitions in \field{d\_hcpcs}.
\end{itemize}

\paragraph{Measurements.}
\begin{itemize}[itemsep=2pt, parsep=1pt]
	\item \field{labevents}: laboratory results (chemistry, hematology, blood gases, etc.)
	      with reference ranges and flags.  Concept definitions in \field{d\_labitems}.
	      Multiple measurements may share a \field{specimen\_id} (same blood draw).
	\item \field{microbiologyevents}: microbiology cultures with specimen type, organism,
	      antibiotic sensitivities, and dilution results.  Stored as a directed hierarchy
	      (specimen $\to$ organism $\to$ antibiotic).
	\item \field{omr}: Online Medical Record outpatient measurements - blood pressure, height,
	      weight, BMI, eGFR - available from both inpatient and outpatient visits.
\end{itemize}

\paragraph{Medications.}
Four tables track prescription and administration:
\begin{itemize}[itemsep=2pt, parsep=1pt]
	\item \field{prescriptions} + \field{pharmacy}: provider orders with drug name, dose,
	      route, and formulary details.
	\item \field{emar} + \field{emar\_detail}: electronic Medicine Administration Record
	      (barcode-scanned actual administrations).  One row per administration event,
	      linked to \field{pharmacy} and \field{poe} tables.
\end{itemize}

\paragraph{Administration.}
\field{services} (admitting service), \field{poe} and \field{poe\_detail} (provider order
entry - orders for imaging, consults, nutrition, medications).

\bigskip
\subsection{ICU Module}
\label{sec:icu}

The \textbf{icu} module (94,458 ICU stays, 65,366 unique patients) contains bedside data
from the MetaVision clinical information system (iMDsoft), limited to patients admitted to
an ICU.  All events reference a \field{stay\_id} (unique per ICU stay), defined in
\field{icustays} with admission/discharge times and care unit names (e.g.\ MICU, SICU, CCU).

\paragraph{Event tables.}
Each documented item is an ``event'' grouped by data type:

\begin{center}
	\small
	\begin{tabular}{lll}
		\tr
		\textbf{Table}           & \textbf{Content}                                     & \textbf{Data Type}       \\
		\mr
		\field{chartevents}      & Vital signs, GCS, vent settings, nursing assessments & Numeric time-series      \\
		\field{inputevents}      & IV infusion rates/amounts (vasopressors, sedation)   & Numeric                  \\
		\field{outputevents}     & Urine output, drain volumes                          & Numeric                  \\
		\field{procedureevents}  & Ventilation, CRRT, catheters (start/end times)       & Categorical + timestamps \\
		\field{datetimeevents}   & Date-valued events (e.g., last dialysis)             & Dates                    \\
		\field{ingredientevents} & Ingredient-level breakdown of infusions              & Numeric                  \\
		\field{d\_items}         & Dimension table for \field{itemid}                   & Reference                \\
		\bottomrule
	\end{tabular}
\end{center}

\paragraph{Suitability for embeddings.}
ICU data is predominantly \textbf{numerical time-series} (vital signs, infusion rates, urine output) rather than natural language.  These signals are not directly embeddable with sentence transformers, and converting them to text (e.g., ``heart rate 92 bpm at 14:30'') would produce low-information chunks unlikely to match condition-level queries semantically.  The ICU module is not a primary source for our corpus.
\paragraph{Metadata enrichment.}
ICU data could strengthen the \textbf{contextual embedding prefix}
(Section~\ref{sec:contextual}) for patients who had ICU stays.  The
\texttt{mimic-code} repository already provides SQL concepts for severity
scores (SOFA, SAPS-II, OASIS), ventilation status, and vasopressor use.
A richer prefix:
\begin{verbatim}
  Discharge note for elderly male, admitted for sepsis
  (SOFA 14, mechanically ventilated, on vasopressors).
  Comorbidities: CHF, CKD.
  Section: Brief Hospital Course.
\end{verbatim}
is more informative than one without ICU severity signals.  However, this enrichment
only matters for \textbf{ICU-relevant conditions} (sepsis, respiratory failure, shock,
post-cardiac arrest) where severity scores carry discriminative information.  For routine
admissions (e.g.\ elective surgery, uncomplicated pneumonia) it maybe doesn't make sense.
%% ============================================================
\part{QA Dataset Generation Pipeline}
\label{sec:architecture}

\section{Pipeline draft}
\subsection{Phases}
\phase{Phase 1}{Corpus Construction}

\begin{enumerate}[leftmargin=*, itemsep=2mm]
	\item \textbf{Select target conditions.}  Query \field{diagnoses\_icd}
	      for ICD code groups with $\geq 200$ (TBD based on dataset statistics. TODO: gather analytics) admissions.\newline
	      Candidate conditions: stroke, myocardial infarction, pneumonia, sepsis, heart failure, COPD exacerbation, acute kidney injury, pulmonary embolism, diabetic ketoacidosis, cirrhosis, etc.  Target: 50--100 conditions.\newline
	      Prefer conditions with high \textbf{ICD co-occurrence richness}: conditions whose admissions carry diverse comorbidity profiles (e.g., pneumonia + CKD, pneumonia + CHF, pneumonia + diabetes) are more likely to produce multi-aspect candidate pools than conditions with uniform ICD profiles. This can be measured as the number of distinct ICD code sets across admissions for a given condition.

	\item \textbf{Build per-condition candidate pools.}  For each condition, gather all discharge note chunks from admissions with that ICD code. This is the full candidate set $D$ for facility-location.

	\item \textbf{Chunk by section + sub-problem.}
	      \begin{itemize}[itemsep=2pt, parsep=1pt]
		      \item Split on section headers
		      \item Within \texttt{Brief Hospital Course}, split on \texttt{\#}~problem markers
	      \end{itemize}

	\item \textbf{Attach metadata to each chunk:}
	      \begin{itemize}[itemsep=2pt, parsep=1pt]
		      \item \field{subject\_id}, \field{hadm\_id}, section name
		      \item All ICD codes for that admission
		      \item Demographics: age group, gender, race, insurance
		      \item Admission type, discharge disposition
	      \end{itemize}

	\item \textbf{Deduplicate same-patient chunks.}
	      The same patient can have multiple hospitalizations with nearly identical boilerplate sections, e.g., \texttt{Past Medical History}, \texttt{Allergies}, and \texttt{Family History} are often copied verbatim across admissions.

	      \textbf{Strategy}: compute a content hash (e.g., SHA-256 of normalized text) per chunk and deduplicate by (\field{subject\_id},\, \text{section},\, \text{hash}). Check also if it's worth considering \emph{only} the most up to date note for a certain patient for a subset of repetitive sections (Past Medical History, etc.)
\end{enumerate}

\phase{Phase 2}{Embeddings Generation}
\label{sec:contextual}

Raw chunks are too patient-specific to match condition-level queries. Following the \textbf{Contextual Retrieval} approach~\cite{anthropic2024contextual}, prepend a metadata prefix to each chunk before embedding:

\begin{quote}
	\small
	\texttt{[age\_group] [gender] [race] patient admitted for [primary ICD description].
	\newline Chief complaint: [chief\_complaint].
	\newline Comorbidities: [top ICD descriptions].
	\newline Section: [section\_name].}
\end{quote}

Check out Section~\ref{sec:embeddings-sec} for more details.

The specific contents to include are to be decided, based on the queries we want to build.

\phase{Phase 3}{Query Generation}
\label{sec:querygen}

\begin{enumerate}[leftmargin=*, itemsep=2mm]
	\item Pick a primary condition (e.g., stroke, ICD-10: I63.*)
	\item Pick 1--2 \textbf{modifier axes} from metadata. Each axis resolves to a filter on \field{hadm\_id}, which is the join key to the discharge notes already chunked in Phase~1. Since chunk metadata was pre-attached in Phase~1 step~4, query generation simply selects which modifier values to combine into the prompt template.

	      \begin{itemize}[itemsep=2pt, parsep=1pt]
		      \item \textbf{Comorbidity} (prior MI, diabetes, renal failure): co-occurring ICD codes on the same admission. A single \field{hadm\_id} carries 5--20 ICD rows in \importantfield{diagnoses\_icd}; ``stroke + diabetes'' is the set of \field{hadm\_id}s that have both \icd{I63.*} and \icd{E11.*}. The human-readable description for the prompt comes from \importantfield{d\_icd\_diagnoses}.

		      \item \textbf{Demographic} (elderly, women, specific race): requires joining two tables. \importantfield{patients} provides \field{gender} and \field{anchor\_age}; \importantfield{admissions} provides \field{race}, \field{insurance}, \field{marital\_status}, and \field{admission\_type}. Both join to \field{diagnoses\_icd} via \field{subject\_id} and \field{hadm\_id} respectively.

		      \item \textbf{Complication} (post-surgical, recurrent, hemorrhagic transformation): draws from multiple sources. Co-occurring ICD codes capture complications (e.g., stroke \icd{I63.*} + hemorrhage \icd{I61.*} on the same \field{hadm\_id}). \field{procedures\_icd} identifies post-surgical modifiers. Recurrence is detected by finding the same \field{subject\_id} with the target condition across multiple \field{hadm\_id}s.
	      \end{itemize}
	\item Prompt an LLM:
	      \begin{quote}
		      \small
		      \emph{``Generate a clinical management question about [condition] that requires knowing how treatment or workup changes when [modifier] is present.  The answer should require synthesizing management decisions from multiple patient cases, not textbook symptom lists.''}
	      \end{quote}
	\item \textbf{Divergence pre-filter.} Before investing in expensive gold annotation, run facility-location and top-$k$ on the candidate pool for the generated query and compute two cheap signals:
	      \begin{itemize}[itemsep=2pt, parsep=1pt]
		      \item \textbf{Jaccard divergence}: if $J(S_\text{cov}, S_\text{topk}) \approx 1$, coverage selects nearly the same chunks as top-$k$ $\Rightarrow$ skip this query, coverage adds nothing here.  If $J \ll 1$, coverage diverges $\Rightarrow$ proceed.
		      \item \textbf{Facility-location gap}: $g(S_\text{cov}) - g(S_\text{topk})$ quantifies how much pool structure top-$k$ ignores. Both $S_\text{cov}$ and $S_\text{topk}$ are already computed, so this is free.
	      \end{itemize}
	      Only queries where coverage meaningfully diverges from top-$k$ proceed to gold annotation.

	\item \textbf{Gold answer generation}

	      The idea: give an LLM access to the candidate pool, have it answer the query comprehensively, and for each piece of information in the answer \textbf{cite which chunk(s)} it drew from.  The cited chunks become the \textbf{gold set}, each tagged with a facet label.

	      \paragraph{Facet annotation.}
	      The gold answer naturally covers several distinct information aspects --- each one is a \textbf{facet}. Facet labels are not pre-defined; the LLM generates them during annotation. For each piece of information in its answer, the LLM cites which chunk(s) it drew from and assigns a short descriptive label for the aspect of care it addresses.

	      \smallskip
	      \textbf{Example.} For the query \emph{``How does management of acute stroke differ when the patient has concurrent CKD?''}, the annotator might produce:

	      \smallskip
	      \small
	      \begin{tabular}{@{}p{0.2\linewidth}p{0.35\linewidth}p{0.37\linewidth}@{}}
		      \textbf{Facet label}          & \textbf{What it captures}                     & \textbf{Example chunk content}                       \\
		      \hline
		      \texttt{anticoag\_adjustment} & CKD affects blood thinner dosing              & Heparin dose reduced to 10u/kg/hr due to CrCl $< 30$ \\[3pt]
		      \texttt{contrast\_avoidance}  & CKD contraindicates contrast imaging          & CT angiography deferred, nephropathy risk            \\[3pt]
		      \texttt{bp\_target}           & Different BP targets in stroke + CKD          & Permissive HTN limit adjusted for renal function     \\[3pt]
		      \texttt{dialysis\_timing}     & When to initiate dialysis during acute stroke & Nephrology consulted for HD timing peri-stroke       \\
	      \end{tabular}
	      \normalsize
	      \smallskip

	      Multiple chunks can map to the same facet (e.g., two different patients' notes both discuss anticoagulation adjustment). Aspect recall does not require retrieving \emph{all} chunks for a facet --- just \emph{at least one} per facet (the $\exists$ in the AR formula, Section~\ref{sec:evaluation}).

	      This is where the coverage hypothesis applies: top-$k$ might retrieve several chunks from the dominant facet (\texttt{anticoag\_adjustment}, which appears in many notes) while missing niche facets (\texttt{dialysis\_timing}, which appears in fewer notes and gets buried by similarity ranking). Coverage should recover more facets by pulling representatives from different regions of the pool.

	      \paragraph{Exhaustive tagging and strict evaluation.}
	      The map pass reads \emph{every} chunk in the pool, so the LLM is instructed to tag \textbf{all} chunks that support each facet, not just the single best one. This makes the gold set comprehensive. At evaluation time, aspect recall uses \textbf{strict chunk-id matching}: a retrieved chunk contributes to a facet only if it appears in the gold set for that facet. No LLM judge is needed at evaluation time --- the check is pure set intersection, making evaluation deterministic and reproducible.

	      \paragraph{Gold annotation via map-reduce.}
	      Since the corpus is heavily pruned before this stage (only conditions with the relevant ICD codes, only queries that pass the divergence filter), each per-condition candidate pool should not be too large. However, to face the long-context that the candidates may give rise to, we can:

	      \begin{enumerate}[itemsep=2pt, parsep=1pt]
		      \item \textbf{Map}: split the pool into batches of $\sim$50 chunks ($\sim$10K tokens).  For each batch, the LLM extracts facts relevant to the query, cites which chunk(s) each fact came from, and assigns a facet label to each citation.
		      \item \textbf{Reduce}: a merge pass synthesizes across batches, deduplicates facets, and produces the final gold answer with the complete (facet $\to$ chunk-ids) mapping.
	      \end{enumerate}

	\item \textbf{(Optional) Query quality verifier}
	      \begin{itemize}[itemsep=2pt, parsep=1pt]
		      \item A separate reader-LLM attempts to answer from basline top-$k$ chunks
		      \item If it succeeds $\Rightarrow$ query does \emph{not} require
		            diversity $\Rightarrow$ discard
		      \item If it fails but succeeds with gold chunks $\Rightarrow$
		            query requires diversity $\Rightarrow$ keep
	      \end{itemize}

\end{enumerate}

\phase{Phase 4}{Evaluation}
\label{sec:evaluation}

\begin{itemize}[itemsep=2pt, parsep=1pt]
	\item \textbf{Aspect recall}: fraction of annotated
	      facets covered by selected $k$ chunks (strict chunk-id matching).\newline
	      $\text{AR}(S) = |\{f \in F : S \cap G_f \neq \emptyset\}| / |F|$, where $G_f$ is the set of gold chunk-ids tagged with facet $f$.
	\item \textbf{Fac score}: facility-location objective value
	      $g(S) = \sum_{i \in D} \max_{j \in S} w_{ij}$
	\item \textbf{Jaccard}: overlap between method's selection and top-$k$ baseline
	\item \textbf{DocRec / FactRec}: gold document/fact fraction recovered
\end{itemize}

\textbf{Methods to compare}: top-$k$ (baseline), MMR, gMMR,
facility-location

\subsection{Does Coverage Pull in Noise?}
\label{sec:coverage-noise}

Facility-location maximizes $g(S) = \sum_{i \in D} \max_{j \in S} w_{ij}$ over the entire candidate pool $D$.  It tries to \emph{cover all of $D$}, including data points that are close to the query in embedding space but irrelevant to the query's specific aspects.  For example, a query about ``pneumonia management in CKD (Complex Kidney Disease) patients'' retrieves a candidate pool where most chunks are about pneumonia, some about CKD, and some are noise (superficially similar embeddings about unrelated respiratory complaints).  Pure coverage would ``waste'' a selection on representing the noise region.

Mitigations:
\begin{enumerate}[leftmargin=*]
	\item \textbf{Candidate pool quality.}  Coverage is applied as a \emph{re-ranker} over $N$ candidates retrieved by ANN (Section~\ref{sec:computational}), not over the full corpus.  If ANN retrieval is reasonable, most of $D$ is query-relevant and noise is a small fraction with weak pull on the objective.
	\item \textbf{$\lambda$ tuning.}  At $\lambda = 1$, coverage degrades to top-$k$; at $\lambda = 0$, it ignores relevance entirely.  The optimal $\lambda$ is a thesis parameter.  If noise is a problem, it manifests as coverage underperforming top-$k$ at low $\lambda$, which we can detect in the experiments.
\end{enumerate}

\subsection{Computational Feasibility}
\label{sec:computational}

Greedy facility-location is $O(|D| \cdot k)$ per selection step, so $O(|D| \cdot k^2)$ total.  With $|D| \approx 3$M chunks (331K notes $\times$ $\sim$10 chunks each) and $k = 10$, this is $\sim 3 \times 10^8$ similarity computations per query, impractical.

Proposal: \textbf{two-stage pipeline}:
\begin{enumerate}[leftmargin=*]
	\item \textbf{Retrieval}: HNSW index retrieves $N \gg k$ candidates by approximate nearest neighbor search.  Typical $N$: 100--1000.  This is sub-linear in $|D|$.
	\item \textbf{Re-ranking}: run facility-location (or MMR, gMMR, FPS) on the $N$ candidates to select $k$.
\end{enumerate}

\paragraph{$N$ as a parameter.} Too small: coverage has no room to diversify (all $N$ candidates are near-duplicates from the dominant cluster).  Too large: distant, irrelevant chunks enter the pool.  $N$ should be swept in experiments alongside $\lambda$ and $k$.

%% ============================================================
\section{LLM Usage and Data Use Agreement}
\label{sec:llm-dua}

Phases~3--4 require feeding MIMIC discharge note chunks into an LLM (query generation, gold answer annotation, optional quality verification).  The \href{https://physionet.org/content/mimiciv/view-license/3.1/}{PhysioNet Credentialed Health Data License~1.5.0} and \href{https://physionet.org/content/mimiciv/view-dua/3.1/}{DUA~1.5.0} impose two constraints:

\begin{enumerate}[leftmargin=*]
	\item \textbf{No sharing.} ``\emph{I will not share access to PhysioNet restricted data with anyone else.}''  Sending clinical text to a commercial API arguably violates this clause.
	\item \textbf{Security.} ``\emph{I will exercise all reasonable and prudent care to maintain the physical and electronic security of PhysioNet restricted data.}''  Even with zero-retention API tiers, the data still transits and is processed on third-party servers outside our control.
\end{enumerate}

\paragraph{Local inference.}
All LLM steps that touch patient data should run on a self-hosted model.  Available hardware: NVIDIA L4 (23\,GB VRAM). Speed? Performance of small LLMs?

\paragraph{Exception.}
Steps that use only \emph{public} information (ICD code descriptions, condition names, modifier labels) can use commercial APIs, since no MIMIC text is transmitted.  This applies to query template generation (Phase~3, step~3) but not to gold annotation.

%% ============================================================
\section{Embeddings: considerations}
\label{sec:embeddings-sec}

\subsection{Contextual Embeddings}

A raw chunk from Brief Hospital Course might read:
\begin{quote}
	\small
	\emph{``Pt had L-sided weakness, slurred speech, presented within 3h
		window.  tPA administered.  CT head negative for hemorrhage.''}
\end{quote}

Without context, embedding this chunk places it in a generic ``neurological emergency'' region.  A query about ``stroke symptoms'' may not rank it highly, because the semantics is different.

\paragraph{Contextual Retrieval~\cite{anthropic2024contextual}}
Prepend contextual information to each chunk before embedding. Reported: 49\% reduction in retrieval failure rate (combined with BM25 hybrid).\\
The original method uses an LLM to generate the prefix from the surrounding document, since document structure is unknown in the general case.  MIMIC discharge notes, however, come with structured metadata (ICD codes, demographics, admission details) already parsed in relational tables, so we replace the LLM call with a deterministic template.

% \paragraph{Alternative: hierarchical two-stage retrieval.}
% Instead of concatenating prefix + chunk into one embedding, decouple them:
% \begin{enumerate}[itemsep=2pt, parsep=1pt]
% 	\item \textbf{Stage 1 (document-level):} embed only the contextual prefix (condition, demographics, chief complaint).  Retrieve relevant \emph{documents} (patients).
% 	\item \textbf{Stage 2 (chunk-level):} within selected documents, embed and retrieve individual chunks using raw clinical text.
% \end{enumerate}
% Stage~1 is cheap (one vector per document) and filters by patient profile; stage~2 is precise and finds the specific section answering the query.  Coverage could be applied at either stage: diverse patients (stage~1) or diverse information pieces (stage~2).  However, this adds a design dimension (where to apply diversity?) and complicates the controlled comparison, which is the thesis's core contribution.  Start with single-stage contextual embeddings; hierarchical retrieval is a second-phase experiment if the single-stage results are inconclusive.

\subsection{Chunking vs.\ Whole-Document Embedding}

Long-context embedding models (e.g., BioClinical ModernBERT with 8K tokens, or Jina v2 with 8K tokens) can embed an entire discharge note as a single vector.  Should we consider?

A whole-document embedding of a 2{,}000-token discharge note compresses heterogeneous content (chief complaint, clinical reasoning,...) into a single 768-dim vector.  Mean-pooling over these diverse sections produces a vector that sits in a ``generic admission'' region of the embedding space, rather than near any specific clinical concept.  This flattens the multi-cluster structure we need, so coverage brings no advantage.

\paragraph{When whole-document embedding?}
If queries target \emph{patient-level} distinctions rather than \emph{information-level} distinctions (e.g., ``find patients similar to this case'' rather than ``what are the symptoms of X and how do they differ with Y''), whole-document embedding is more natural.  This is based upon the queries we want to build, so maybe it's good

\subsection{Embedding Models for medical domain}
\label{sec:embeddings}

\begin{itemize}[itemsep=4pt, parsep=1pt]

	\item \textbf{BioLORD-2023}~\cite{biolord}: sentence-transformer fine-tuned from \texttt{all-mpnet-base-v2}. 768-dim embeddings, 512-token context.

	\item \textbf{MedEmbed}~\cite{medembed}: family of sentence-transformers (small/base/large) fine-tuned from the BAAI/BGE series. The large variant produces 1024-dim embeddings, 512-token context.  Apache-2.0 licence.
	      \
	\item \textbf{BioClinical ModernBERT Embeddings~\cite{sounack2025bioclinical}}: sentence-transformer fine-tuned from BioClinical ModernBERT, a ModernBERT variant further pre-trained on PubMed abstracts and clinical text in the MIMIC lineage.  768-dim embeddings, up to 8{,}192-token context. Apache-2.0 licence.

\end{itemize}

%% ============================================================
\section{Cluster Validation}
\label{sec:validation}

Do candidate pools have the \hyperref[sec:introduction]{desired multi-cluster structure of Section~\ref*{sec:introduction}}?

\subsection{Pipeline draft}

\begin{enumerate}[leftmargin=*]
	\item \textbf{Embed} all chunks with \texttt{multi-qa-mpnet-base-cos-v1}
	      (768 dimensions).

	\item \textbf{UMAP} reduction to $\sim$15--30 dimensions.  UMAP~\cite{mcinnes2018umap} preserves local \emph{and} global structure better than t-SNE.

	\item \textbf{HDBSCAN}~\cite{campello2013hdbscan} with \field{min\_cluster\_size} tuned per condition.

	\item \textbf{Validate} structure with:
	      \begin{itemize}[itemsep=2pt, parsep=1pt]
		      \item \textbf{DBCV}~\cite{moulavi2014dbcv}: density-based cluster
		            validation, designed for HDBSCAN output
		      \item \textbf{Cluster size ratio}: $|C_{\max}| / |C_{\text{total}}|$
		            measures dominance of central cluster
		      \item \textbf{Density ratio}: $\rho_{\text{central}} / \rho_{\text{niche}}$
		      \item \textbf{Number of clusters + noise fraction}
	      \end{itemize}

	\item \textbf{Cross-reference} HDBSCAN clusters with ICD sub-codes.  If embedding clusters align with clinical sub-types, the structure is semantically meaningful.
\end{enumerate}

\subsection{Is Full Clustering Necessary?}
\label{sec:simpler-validation}
The Jaccard divergence and facility-location gap signals used in Phase~3 step~4 already tell us whether coverage selects differently from top-$k$, without requiring any clustering. Full clustering (UMAP + HDBSCAN) is not needed as a pipeline gate, but remains useful as a post-hoc \emph{explanatory} analysis: correlating cluster statistics (number of clusters, dominance ratio) with $\Delta$AR can help explain \emph{why} coverage helped on certain conditions. Ultimately, \textbf{aspect recall} is the real validation: if coverage achieves higher aspect recall than top-$k$ on a condition, the pool structure was there.

\subsection{High-Dimensional Clustering: Caveats}
\label{sec:clustering-caveats}

If we do use clustering for analysis, the following caveats apply.

\paragraph{Curse of dimensionality.}
HDBSCAN relies on point density.  In high dimensions, pairwise distances concentrate~\cite{aggarwal2001surprising} and density-based clustering becomes unreliable.  This motivates the UMAP$\,+\,$HDBSCAN pipeline (used by BERTopic~\cite{grootendorst2022bertopic}, Top2Vec~\cite{angelov2020top2vec}).

\paragraph{Is UMAP necessary?}
\label{sec:umap-necessity}
The concentration-of-distances argument applies to \emph{uniformly distributed} high-dimensional data.  Sentence embeddings are not uniform, they occupy a narrow anisotropic cone~\cite{ethayarajh2019contextual}, and their intrinsic dimensionality is much lower than the ambient 768.  HDBSCAN with cosine distance may find meaningful clusters directly in the original space \emph{if} the effective dimensionality is low enough.  This is an empirical question: run HDBSCAN both with and without UMAP and compare.  If results agree, UMAP adds stochastic noise (sensitive to \field{n\_neighbors}, \field{min\_dist}) without benefit.

\paragraph{Anisotropy cuts both ways.}
Anisotropy means clusters may reflect frequency or style artifacts (all discharge notes share boilerplate language) rather than the semantic distinctions we care about.  Cross-referencing HDBSCAN clusters with ICD sub-codes (Step~5 above) is essential: if clusters don't align with clinical sub-types, they are artifacts.

\bibliographystyle{plain}
\bibliography{mimic-feasibility-analysis}

\appendix
\section*{Appendix}
\addcontentsline{toc}{section}{Appendix}

\section{Available Analytics from \texttt{mimic-code}}
\label{sec:mimic-code}

The official \texttt{mimic-code} repository provides SQL concepts, cohort definitions, and mapping tables that we can use.

\begin{center}
	\small
	\begin{tabular}{p{4.2cm}p{5.4cm}p{6cm}}
		Resource                                                       & Location & Role in Our Pipeline \\
		\mr
		\textbf{Charlson Comorbidity Index}                            &
		\field{mimic-iv/concepts/\allowbreak comorbidity/charlson.sql} &
		Maps ICD-9/10 codes to 16 clinical categories (MI, CHF, stroke, COPD,
		diabetes, renal, cancer\ldots).  Primary tool for \textbf{condition
		selection} and \textbf{comorbidity modifier axes}.                                               \\
		\addlinespace
		\textbf{CCS hierarchical grouping}                             &
		\field{mimic-iii/concepts\_postgres/\allowbreak diagnosis/ccs\_dx.sql}
		+ \field{ccs\_multi\_dx.csv.gz}                                &
		4-level hierarchical ICD-9 $\to$ clinical category mapping.
		Useful for building the macrocategory index in hierarchical retrieval
		(Phase~2).                                                                                       \\
		\addlinespace
		\textbf{Concept maps}                                          &
		\field{mimic-iv/concepts/\allowbreak concept\_map/}            &
		ICD $\to$ LOINC/OMOP/SNOMED CSVs.  Cross-referencing structured data
		when building contextual embedding prefixes.                                                     \\
		\addlinespace
		\textbf{Age / demographics}                                    &
		\field{mimic-iv/concepts/\allowbreak demographics/age.sql}     &
		Age cohort extraction.  Needed for demographic modifier axes.                                    \\
		\addlinespace
		\textbf{Severity scores}                                       &
		\field{mimic-iv/concepts/\allowbreak score/} (SOFA, SAPS-II, LODS,
		APS-III, OASIS)                                                &
		Patient acuity signals.  Could serve as additional metadata for
		contextual embeddings or as a proxy for clinical complexity.                                     \\
	\end{tabular}
\end{center}

\end{document}

